\subsection{Cronograma de Sprints y Fases del Proyecto}

El desarrollo se ha estructurado en cinco \textit{Sprints}, integrando no solo la implementación técnica, sino también la labor de investigación y análisis necesaria para el rigor del Trabajo de Fin de Grado.

\begin{itemize}
    \item \textbf{Sprint 1: Análisis e Investigación (Fase Previa)}:
    Dedicado al estudio del estado del arte y fundamentos teóricos.
    \begin{itemize}
        \item \textbf{Marco Socioeconómico}: Análisis de la necesidad de plataformas de intercambio lingüístico y el impacto de la globalización en el aprendizaje de idiomas.
        \item \textbf{Marco Regulatorio}: Estudio del cumplimiento de la normativa \textbf{RGPD} (Reglamento General de Protección de Datos) para el manejo de correos y datos de usuarios.
        \item \textbf{Investigación Tecnológica}: Selección del \textit{stack} tecnológico (MERN vs. otros) y evaluación de la viabilidad de Next.js para los requerimientos del sistema.
    \end{itemize}

    \item \textbf{Sprint 2: Diseño de Arquitectura y Autenticación}:
    Definición del modelo de datos en MongoDB, esquemas de Mongoose y desarrollo del sistema de seguridad basado en JWT y cookies \textit{HttpOnly}.

    \item \textbf{Sprint 3: Desarrollo del Core y Lógica de Negocio}:
    Implementación del sistema de creación de cartas, flujo de envío entre usuarios y lógica de corrección.

    \item \textbf{Sprint 4: Funcionalidades Sociales y Multimedia}:
    Desarrollo del sistema de seguidores (\textit{Follows}), solicitudes de amistad y gestión de archivos estáticos con Multer para las imágenes de perfil.

    \item \textbf{Sprint 5: Pruebas, Refactorización y Conclusiones}:
    Fase final de depuración de errores (\textit{debugging}), optimización del rendimiento y redacción de las conclusiones finales de la memoria.
\end{itemize}